%% ============================================================
%% CHAPTER 5 — RESULTS AND INTEGRATED ANALYSIS
%% ============================================================

\chapter{Results and Integrated Analysis}
\label{ch:results}

\section{Chapter Overview}
\label{sec:results_overview}

This chapter presents the unified results of the four-layer framework,
integrating the physics engine (Zone~C) and the adversarial NDE model
(Zone~A) into a coherent analytical picture.  \Cref{sec:results_network}
examines the network vulnerability landscape.
\Cref{sec:results_stackelberg} quantifies the benefit of physics-informed
over uniform defender allocation.  \Cref{sec:results_adv_network} links
the adversarial ASR to the game-theoretic attacker strategy.
\Cref{sec:results_dashboard} describes the interactive dashboard as a
decision-support tool.  \Cref{sec:results_discussion} synthesises the
three research questions posed in \cref{sec:intro_problem}.

\section{Network Vulnerability Landscape}
\label{sec:results_network}

\Cref{fig:network_map} in \cref{ch:physics} revealed a strongly
heterogeneous $\Pf$ landscape: 8 of 22 segments carry $\Pf > \num{0.75}$,
yet these segments carry \SI{64}{\percent} of the total normalised
network value (diameter $\times$ length product).  The Pearson correlation
between $\Pf$ and betweenness centrality is $\rho = \num{+0.41}$,
indicating a moderate but significant tendency for physically vulnerable
segments to also lie on critical routing paths --- the precise combination
that a state-actor adversary (\Cref{tab:attacker_profiles}) is designed
to exploit.

\begin{figure}[htbp]
  \centering
  \includegraphics[width=0.92\textwidth]{fig12_network_property_distributions}
  \caption{Four-panel network property analysis.
    \textit{Top-left:} Diameter distribution.
    \textit{Top-right:} $\Pf$ distribution across 22 segments.
    \textit{Bottom-left:} Betweenness centrality.
    \textit{Bottom-right:} Scatter plot of $\Pf$ versus normalised
    segment value, coloured by seam type.  The positive correlation
    between structural vulnerability and network centrality provides
    the strategic rationale for physics-informed resource allocation.}
  \label{fig:network_props}
\end{figure}

\section{Benefit of Physics-Informed Stackelberg Allocation}
\label{sec:results_stackelberg}

\subsection{Scenario Definitions}
\label{subsec:results_scenarios}

Three defender scenarios are compared:

\begin{enumerate}
  \item \textbf{Unprotected}: $c_i = 0$ for all $i$. Represents the
    pre-implementation baseline.
  \item \textbf{Uniform baseline}: $c_i = B/n = \num{0.30}$ for all $i$.
    Represents naive equal-allocation without physics information.
  \item \textbf{Physics-informed SSE}: $\bmc^*$ solving the Bayesian
    SSE as formulated in \cref{subsec:physics_lp}.
\end{enumerate}

\subsection{Total Expected Network Loss}
\label{subsec:results_loss}

Total expected network loss under coverage $\bmc$ and protection
factor $\delta$ is:
\begin{equation}
  \mathcal{L}(\bmc) = \sum_{e \in \mathcal{E}}
    \Pf(e) \cdot v(e) \cdot \bigl[1 - c_e(1 - \delta)\bigr],
  \label{eq:network_loss}
\end{equation}
where $v(e)$ is the normalised segment value and $\delta = \num{0.25}$.
\Cref{tab:scenario_comparison} presents the results.

\begin{table}[htbp]
  \centering
  \caption{Scenario comparison: total expected network loss under three
    defender strategies.  All values computed on the 22-segment Gulf
    Coast network with $\delta = \num{0.25}$ and $B = \num{6.6}$.}
  \label{tab:scenario_comparison}
  \begin{tabular}{@{}lcccc@{}}
    \toprule
    Scenario & $\mathcal{L}(\bmc)$ & vs.\ Unprotected & vs.\ Baseline \\
    \midrule
    Unprotected          & 3.830 & ---             & --- \\
    Uniform Baseline     & 2.968 & $-\SI{22.5}{\percent}$ & --- \\
    Physics-Informed SSE & 2.465 & $-\SI{35.7}{\percent}$
      & $-\SI{17.0}{\percent}$ \\
    \bottomrule
  \end{tabular}
\end{table}

The key finding: physics-informed allocation reduces network loss by
\SI{17.0}{\percent} relative to the uniform baseline, at the same
defender budget.  This additional reduction is attributable solely to
the physics-informed payoff function: the SSE concentrates resources
on the subset of segments where $\Pf$ is both high and correlated with
network centrality, rather than spreading them uniformly.

\begin{figure}[htbp]
  \centering
  \includegraphics[width=0.92\textwidth]{fig20_scenario_comparison}
  \caption{Per-segment residual risk under each scenario (upper panel)
    and adversarial robustness curves (lower panel).
    \textit{Upper:} The SSE allocation (navy blue) systematically
    outperforms the uniform baseline (amber) on the high-$\Pf$
    segments (right half of $x$-axis), which carry disproportionate
    network risk.
    \textit{Lower:} Accuracy versus $\eps$ for FGSM, BIM, and PGD.}
  \label{fig:scenario_comparison}
\end{figure}

\section{Adversarial NDE and the Game-Theoretic Attacker}
\label{sec:results_adv_network}

The Bayesian SSE attacker strategy assigns $q^* = $ SEG\_N012\_N015
as the dominant target (probability $\approx \num{0.60}$ for the
\textsc{State Actor} type, \num{0.45} for \textsc{Strategic}).  This
segment has seam type LF-ERW with \acs{SCF} $= \num{1.8}$, yielding
an effective adversarial budget of:
\begin{equation*}
  \eps_\mathrm{eff} = \num{0.30} \times \frac{1.8}{1.5} = \num{0.36}
\end{equation*}
At this $\eps$, the \acs{FGSM} \acs{ASR} is approximately
\SI{24}{\percent}.  In operational terms: if a state-actor adversary
targets SEG\_N012\_N015 and simultaneously manipulates its \ac{NDE}
inspection report, there is a $\sim$24\% probability that the defect
is classified as CLEAN, deferring remediation and allowing the physical
attack timeline to proceed undetected.

The \emph{cross-layer threat model} is thus quantified: the attacker
utility from the Stackelberg game (high $\Pf$, high centrality) and the
adversarial NDE ASR are positively correlated across segments, as shown
in \cref{fig:dashboard_intel} (lower right panel).

\section{Interactive Dashboard}
\label{sec:results_dashboard}

\Cref{fig:dashboard_intel} illustrates the three-tab dashboard delivered
as Sprint~5.  The dashboard serves two audiences simultaneously:

\begin{itemize}
  \item \textbf{Integrity engineers}: the click-to-inspect FAD panel
    shows each segment's assessment point against the BS~7910 Option~1
    curve, with nominal and critical flaw assessments pre-computed.
  \item \textbf{Security planners}: the SSE coverage heat-map and
    scenario comparison tab quantify the cost of uniform vs.\
    physics-informed allocation in expected loss units directly.
\end{itemize}

\begin{figure}[htbp]
  \centering
  \includegraphics[width=0.92\textwidth]{fig18_dashboard_network_intelligence}
  \caption{Sprint~5 dashboard: Network Intelligence tab for the highest-risk
    segment (SEG\_N012\_N015, $\Pf = \num{0.93}$, LF-ERW seam).
    \textit{Top:} Pipeline network coloured by $\Pf$.
    \textit{Middle:} BS~7910 FAD assessment (critical flaw outside
    the envelope).
    \textit{Bottom:} FGSM adversarial feature perturbation by group,
    with attack success rate gauge.}
  \label{fig:dashboard_intel}
\end{figure}

\begin{figure}[htbp]
  \centering
  \includegraphics[width=0.92\textwidth]{fig19_stackelberg_coverage_heatmap}
  \caption{Sprint~5 dashboard: Stackelberg Defender tab.
    \textit{Top-left:} SSE coverage allocation (navy, colour depth
    $\propto\Pf$).
    \textit{Top-right:} Uniform baseline (amber, constant $c_i = \num{0.30}$).
    \textit{Bottom-left:} SSE $-$ baseline gain (green = above baseline,
    red = below).
    \textit{Bottom-right:} Budget sensitivity, matching
    \cref{fig:budget_sensitivity}.}
  \label{fig:dashboard_game}
\end{figure}

\section{Discussion}
\label{sec:results_discussion}

\paragraph{Research Question 1} (Does physics-informed payoff produce a
qualitatively different SSE?)\quad Yes.  The uniform baseline assigns
identical protection to a seamless X65 DSAW segment ($\Pf = \num{0.29}$)
and to an LF-ERW legacy segment ($\Pf = \num{0.93}$).  The
physics-informed SSE assigns \SI{0}{\percent} coverage to the former and
\SI{$>$30}{\percent} to the latter, producing a structurally concentrated
allocation that no abstract constant payoff would generate.

\paragraph{Research Question 2} (How does the SSE compare to uniform
allocation?)\quad The SSE achieves a \SI{17.0}{\percent} reduction in
total expected network loss at the same budget.  This is achieved through
a convex redistribution: the top 5 segments (by $\Pf \times$ centrality)
receive $\approx\SI{60}{\percent}$ of the budget.

\paragraph{Research Question 3} (Can an adversary suppress the $\Pf$
signal via NDE manipulation?)\quad Yes, with \ac{BIM} \acs{ASR} of
\SI{20.1}{\percent} at $\eps = \num{0.30}$.  The physics-informed
$\eps$ scaling (\cref{eq:eps_scaling}) reveals that the most
game-theoretically targeted segments are also those most susceptible to
adversarial NDE manipulation, identifying a compounding risk that neither
the structural integrity nor the security game literature has previously
quantified.
